% Part: methods
% Chapter: induction
% Section: strong-induction

\documentclass[../../../include/open-logic-section]{subfiles}

\begin{document}

\olfileid[hi]{mth}{ind}{str}

\olsection{प्रबल आगमन}

ऊपर चर्चित आगमन सिद्धांत में हम $P(0)$ और यह भी सिद्ध करते हैं कि यदि
$P(k)$, तो $P(k+1)$। दूसरे भाग में हम $P(k)$ को सत्य मानते हैं और इस
मान्यता से $P(k+1)$ सिद्ध करते हैं। निस्संदेह तुल्य रूप से हम $P(k-1)$ मानकर
उससे $P(k)$ सिद्ध कर सकते थे---महत्त्वपूर्ण बात यह है कि किसी भी संख्या से
उसके उत्तरवर्ती तक अनुमान कर सकें; कि किसी भी संख्या के लिए संबंधित दावे को
यह मानकर सिद्ध कर सकें कि वह उसके पूर्ववर्ती के लिए सत्य है।

आगमन सिद्धांत का एक रूप ऐसा है जिसमें हम केवल यह नहीं मानते कि दावा $k$ के
पूर्ववर्ती $k-1$ के लिए सत्य है, बल्कि यह मानते हैं कि वह $k$ से छोटी सभी
संख्याओं के लिए सत्य है, और इस मान्यता से $k$ के लिए दावा स्थापित करते हैं।
इससे भी सभी $n\in\Nat$ के लिए दावा $P(n)$ मिलता है। क्योंकि $P(0)$ स्थापित
करते ही हमने स्थापित कर दिया कि $P$, $1$ से छोटी सभी संख्याओं के लिए सत्य
है। और यदि हम जानते हैं कि सभी $l<k$ के लिए $P(l)$ हो तो $P(k)$, तो इसे
विशेषतः $k=1$ के लिए जानते हैं। अतः $P(1)$ निष्कर्षित कर सकते हैं। इससे हमने
$P(0)$ और $P(1)$, अर्थात् सभी $l<2$ के लिए $P(l)$ सिद्ध कर दिया, और चूँकि
हमारे पास यह सशर्त भी है कि यदि सभी $l<2$ के लिए $P(l)$, तो $P(2)$, हम
$P(2)$ निष्कर्षित कर सकते हैं, और इसी प्रकार आगे।

वास्तव में यदि हम सामान्य सशर्त ``सभी $k$ के लिए, यदि सभी $l<k$ के लिए
$P(l)$, तो $P(k)$'' स्थापित कर सकें, तो हमें $P(0)$ अलग से स्थापित नहीं
करना पड़ता, क्योंकि वह इससे निष्पन्न होता है। याद रखें कि ``सभी $l<k$ के
लिए $P(l)$'' जैसा सामान्य दावा तब सत्य है जब कोई $l<k$ न हो। यह रिक्त
परिमाणीकरण की स्थिति है: ``सभी $A$, $B$ हैं'' तब सत्य है जब कोई $A$ न हो;
$\lforall[x][(!A(x)\lif !B(x))]$ तब सत्य है जब कोई $x$, $!A(x)$ को संतुष्ट
न करे। इस स्थिति में औपचारिक रूप ``$\lforall[l][(l<k\lif P(l))]$'' होगा---और
यदि कोई $l<k$ न हो तो वह सत्य है। और यदि $k=0$, तो ठीक यही स्थिति है: कोई
$l<0$ नहीं, अतः ``सभी $l<0$ के लिए $P(0)$'' सत्य है, चाहे $P$ कुछ भी हो।
इस प्रकार ``यदि सभी $l<k$ के लिए $P(l)$, तो $P(k)$'' का प्रमाण अपने-आप
$P(0)$ स्थापित कर देता है।

यह रूप तब उपयोगी है जब $k$ के लिए दावा स्थापित करना केवल $k-1$ के लिए दावे
पर निर्भर नहीं किया जा सकता, बल्कि एक या अधिक $l<k$ के लिए उसके सत्य होने
की मान्यता चाहिए।

\end{document}
